\section{Conclusões}\label{sec:conclusoes}

Este estudo delineou o panorama global das pesquisas dedicadas a Ligas de Alta Entropia Refratárias (RHEAs) para superação das barreiras aerotérmicas do voo hipersônico ao longo de duas décadas de produção científica (2004--2026). A trajetória do campo evidencia um crescimento acelerado e contínuo, com taxa média de expansão composta de 15,8\% ao ano. A transição de uma fase inicial de curiosidade teórica para um domínio de prioridade estratégica em defesa e infraestrutura aeroespacial reflete-se na concentração massiva de mais de 90\% de toda a produção bibliográfica no intervalo pós-2020, culminando no ápice produtivo verificado em 2025.

A geografia da colaboração científica revela uma divisão funcional clara de competências entre os principais polos globais. A Alemanha e instituições parceiras europeias operam como núcleos epistêmicos da termodinâmica fundamental, concentrando-se nos mecanismos atômicos de oxidação, estabilidade de fases e caracterização de danos sob altas temperaturas. Paralelamente, a China assume a liderança incontestável em volume de publicações e desenvolvimento experimental aplicado, impulsionando a metalurgia do pó avançada, síntese por plasma de faísca (SPS) e manufatura aditiva a laser.

Superar a vulnerabilidade histórica à oxidação catastrófica (\emph{pest oxidation}) constitui o eixo de convergência das frentes atuais de pesquisa. A literatura caminha decididamente para a superação do paradigma de ligas monolíticas através da consolidação de revestimentos de silicietos dopados com elementos de terras raras, dotados de mecanismos de autorregeneração vítrea (\emph{self-healing}), e do emprego de arquiteturas híbridas cerâmica-metal. Concomitantemente, a simbiose entre inteligência artificial e termodinâmica computacional (CALPHAD) substitui abordagens tradicionais de tentativa e erro, acelerando o mapeamento de formulações de alta ductilidade a frio e resistência superior a 1500 \textdegree{}C. As ligas de alta entropia consolidam-se, assim, como pilares estruturais indispensáveis para a nova geração de veículos reutilizáveis e sistemas de propulsão hipersônica.
